\begin{helper}
Let $G$ and $G'$ be finite $\Delta$-regular graphs, let $\gamma$ be a real
parameter, and let $\mu$ and $\mu'$ be activation--priority sampling laws for
$G$ and $G'$ with the same parameters \(\Delta\) and \(\gamma\), hence with
the same activation ratio \(p=\gamma/\Delta\) whenever \(\Delta>0\).  Suppose that
$\phi:V(G)\to V(G')$ is injective, $X\subseteq N_G(r)$, the map $\phi$
preserves all adjacencies inside $X\cup\{r\}$.  Suppose moreover that $G'$
has the following split-blocker interface.  For every original outside
incidence $(x,z)$ with $x\in X$, $z\notin X\cup\{r\}$, and $xz\in E(G)$,
there is a replacement vertex $\beta(x,z)\in V(G')$.  The replacement
vertices are distinct and are fresh from the embedded local copy:
$\beta(x,z)\ne\phi(y)$ for every $y\in X\cup\{r\}$.  Also
$\beta(x,z)$ is adjacent to $\phi(x)$, is not adjacent to $\phi(r)$, and is
not adjacent to any $\phi(y)$ with $y\in X\setminus\{x\}$.  Finally, every
neighbor of an embedded vertex $\phi(x)$ with $x\in X$ is either an embedded
local neighbor $\phi(y)$ with $y\in X\cup\{r\}$ and $xy\in E(G)$, or one of
the replacement vertices $\beta(x,z)$ for an original outside incidence
$(x,z)$.  Then
\[
\sum_{S\subseteq V(G)} {\bf 1}_{S\cap X\ne\emptyset}\mu(S)
\le
\sum_{S'\subseteq V(G')} {\bf 1}_{S'\cap \phi(X)\ne\emptyset}\mu'(S').
\]
\end{helper}
\begin{proof}
By \noderef{RandomIndependentSetSampling}, choose underlying
activation--priority laws \(\nu\) and \(\nu'\) whose push-forwards are
\(\mu\) and \(\mu'\).  Thus an outcome for \(G\) is a pair
\((A,\pi)\), an outcome for \(G'\) is a pair \((A',\pi')\), the activated
sets have the Bernoulli product law with parameter \(\gamma/\Delta\), the
priority coordinates are independent uniform coordinates on \([0,1]\), and
the output sets are
\[
I_G(A,\pi)=\{v\in A:\pi(w)<\pi(v)\hbox{ for all }w\in A\cap N_G(v)\}
\]
and the analogous set \(I_{G'}(A',\pi')\).  The same definition says that
\(\mu\) and \(\mu'\) are precisely the push-forward laws of these two output
maps.

Apply \noderef{SamplingNonemptyPushForwardEvent} first to the source law and
the finite vertex set \(X\), and then to the target law and the finite vertex
set \(\phi(X)\).  It gives the two endpoint identities
\[
\operatorname{ofReal}\!\left(\sum_{S\subseteq V(G)}
{\bf 1}_{S\cap X\ne\emptyset}\mu(S)\right)
=
\nu\{(A,\pi):I_G(A,\pi)\cap X\ne\emptyset\},
\]
and
\[
\operatorname{ofReal}\!\left(\sum_{S'\subseteq V(G')}
{\bf 1}_{S'\cap\phi(X)\ne\emptyset}\mu'(S')\right)
=
\nu'\{(A',\pi'):I_{G'}(A',\pi')\cap\phi(X)\ne\emptyset\}.
\]
Therefore it remains to compare these two activation--priority event
probabilities.

The hypotheses of the present helper are the paper-facing split-blocker
interface.  They give finite \(\Delta\)-regular graphs, two sampling laws with
the same \(\Delta\) and \(\gamma\), an injective local copy \(\phi\), the
containment \(X\subseteq N_G(r)\), preservation of all adjacencies in
\(X\cup\{r\}\), and private replacement vertices \(\beta(x,z)\) that are
distinct, fresh from \(\phi(X\cup\{r\})\), adjacent only to their intended
embedded candidate among vertices of \(\phi(X)\), nonadjacent to \(\phi(r)\),
and exhaustive for nonlocal neighbors of embedded \(X\)-vertices.

Invoke \noderef{SamplingSplitOutsideBlockerHybridRepresentation}.  It gives a finite list
\(z_0,\ldots,z_{n-1}\) of the original outside blockers, the candidate sets
\(B_k=\{x\in X:xz_k\in E(G)\}\), one finite exposed-state set \(\Omega\),
nonnegative weights \(w:\Omega\to\mathbb R\), inner stage probabilities
\(P(k,\omega)\), pairwise disjoint source and target cell families covering
the unit-priority supports up to null complements, the endpoint cell
identities
\[
\operatorname{ofReal}(w(\omega)P(0,\omega))
=
\nu(\{(A,\pi):I_G(A,\pi)\cap X\ne\emptyset\}\cap C_\omega),
\]
\[
\operatorname{ofReal}(w(\omega)P(n,\omega))
=
\nu'(\{(A',\pi'):I_{G'}(A',\pi')\cap\phi(X)\ne\emptyset\}\cap C'_\omega),
\]
and the adjacent monotonicity inequalities
\[
P(k,\omega)\le P(k+1,\omega)\qquad(k<n,\ \omega\in\Omega).
\]
These are statement-level conclusions of the hybrid node; the present proof
does not reconstruct its displayed cells, fixed-stage normalizations, or
candidate-priority comparison.

Now apply \noderef{SamplingFiniteConditionalReplacementIteration} to the
finite set \(\Omega\), the nonnegative weight function \(w\), the stage
function \(P\), and the adjacent inequalities just supplied by
\noderef{SamplingSplitOutsideBlockerHybridRepresentation}.  This gives
\[
\sum_{\omega\in\Omega}w(\omega)P(0,\omega)
\le
\sum_{\omega\in\Omega}w(\omega)P(n,\omega).
\]
Because \(w(\omega)\ge0\), \(P(j,\omega)\ge0\), and \(P(j,\omega)=0\) on
zero-weight atoms, the endpoint cell identities can be summed over the finite
atom family.  Applying
\noderef{SamplingFiniteCellEndpointSumDecomposition} to the source endpoint
event, the source cells \(C_\omega\), and the values \(P(0,\omega)\), using
the source disjointness, cover, null-complement, and endpoint cell-mass
clauses from \noderef{SamplingSplitOutsideBlockerHybridRepresentation}, gives
\[
\operatorname{ofReal}\!\left(\sum_{\omega\in\Omega}w(\omega)P(0,\omega)\right)
=
\nu\{(A,\pi):I_G(A,\pi)\cap X\ne\emptyset\}.
\]
The same finite decomposition lemma applied to the target endpoint event, the
target cells \(C'_\omega\), and the values \(P(n,\omega)\), gives
\[
\operatorname{ofReal}\!\left(\sum_{\omega\in\Omega}w(\omega)P(n,\omega)\right)
=
\nu'\{(A',\pi'):I_{G'}(A',\pi')\cap\phi(X)\ne\emptyset\}.
\]
Monotonicity of \(\operatorname{ofReal}\) applied to the weighted-sum
inequality therefore yields
\[
\nu\{(A,\pi):I_G(A,\pi)\cap X\ne\emptyset\}
\le
\nu'\{(A',\pi'):I_{G'}(A',\pi')\cap\phi(X)\ne\emptyset\},
\]
after substituting the hybrid endpoint event identities.

Finally substitute the two push-forward identities from
\noderef{SamplingNonemptyPushForwardEvent}.  The target sum is nonnegative
because \(\mu'\) is a probability mass function in
\noderef{RandomIndependentSetSampling} and every indicator is nonnegative.
Hence the order comparison of the two \(\operatorname{ofReal}\)-images
implies the corresponding real inequality of the two finite push-forward
masses:
\[
\sum_{S\subseteq V(G)} {\bf 1}_{S\cap X\ne\emptyset}\mu(S)
\le
\sum_{S'\subseteq V(G')} {\bf 1}_{S'\cap\phi(X)\ne\emptyset}\mu'(S').
\]
This is the desired event monotonicity.
\end{proof}
